% Generated draft: requires mathematical and source review.
\documentclass[11pt]{article}
\usepackage[margin=25mm]{geometry}
\usepackage{fontspec}
\setmainfont{Latin Modern Roman}
\usepackage{amsmath,amssymb,mathtools,unicode-math}
\setmathfont{Latin Modern Math}
\usepackage{enumitem,needspace,xurl,hyperref}
\setlength{\parindent}{0pt}
\setlength{\parskip}{5pt}
\setlength{\emergencystretch}{3em}
\newcommand{\N}{\mathbb{N}}
\newcommand{\Z}{\mathbb{Z}}
\newcommand{\Q}{\mathbb{Q}}
\newcommand{\R}{\mathbb{R}}
\newcommand{\C}{\mathbb{C}}
\newcommand{\F}{\mathbb{F}}
\newcommand{\A}{\mathbb{A}}
\newcommand{\PP}{\mathbb{P}}
\newcommand{\E}{\mathbb{E}}
\newcommand{\eps}{\varepsilon}
\newcommand{\dd}{\,\mathrm{d}}
\begin{document}
\section*{UnsolvedMath ID: 20002061}
\textbf{Source problem code:} catalogue entry
\subsection{Definitions and mathematical statement}
The source deliberately records only ``something approximating'' Tobias Ekholm's question and says its author probably did not capture the intended statement. It asks whether a Legendrian isotopic to the zero-section Legendrian must have intersecting cotangent projection, with compactness queried parenthetically. Eliashberg's comment says the literal claim follows from generating functions; the stronger workshop question and its nontrivial proof are not recovered.

For the literal model, let
\[
 J^1(M)=\mathbb R_u\times T^*M,\qquad
 \alpha=du-p\,dq,\qquad
 \rho(u,q,p)=(q,p),\qquad
 L_0=\{(0,q,0):q\in M\}.
\]
The question is whether every Legendrian \(L'\) Legendrian-isotopic to \(L_0\) satisfies
\[
 \rho(L')\cap0_M\ne\varnothing.
\]
This is intersection after projection, not necessarily intersection of \(L'\) with \(L_0\) in the jet space.

The aid reports an affirmative answer when \(M\) is closed. Chekanov persistence supplies a quadratic-at-infinity generating family \(F\). In local fiber coordinates \(\xi\), its fiber-critical locus parametrizes
\[
 (q,\xi)\longmapsto(F(q,\xi),q,d_qF(q,\xi)),
 \qquad d_\xi F=0.
\]
Thus zero-cotangent projection occurs exactly at full critical points \(dF=0\); their critical values need not be zero. Quadratic-at-infinity Morse theory forces such a point. If \(F\) is Morse, the reported number of lifted intersection points, counted as points on \(L'\), is at least
\[
 \sum_j b_j(M;\mathbb F_2).
\]
This need not equal the number of distinct projected cotangent points when different lifts project together. The source's informal fiberwise-critical-value/Cerf-diagram definition of a generating function also requires the usual regularity and parametrization hypotheses to define a smooth Legendrian; these are not stated there.

Without closedness, the literal claim fails. Over \(M=\mathbb R\), the strict contact isotopy
\[
 \Phi_t(u,x,p)=(u+tx,x,p+t),\qquad0\leq t\leq1,
\]
sends \(L_0\) to \(j^1(tx)\); at \(t=1\) its cotangent projection has \(p=1\) and misses the zero section. The embedding is proper, so properness alone does not replace closedness. The same formula gives failure over \([0,1]\) if Legendrians with boundary and this boundary behavior are allowed; compactness with boundary is therefore not the closed-base hypothesis.

For a noncompact base, a compactly supported isotopy, or one supported over a compact proper subset of the base, instead leaves the zero-section Legendrian unchanged outside that support. Its projected intersection is then present for that separate geometric reason. These three regimes do not settle an unspecified stronger version of the workshop question.
\subsection{Short English statement}
For the literal question, the supplied aid gives projected zero-section intersection on a closed base via quadratic-at-infinity generating families. A strict jet translation disproves the unrestricted open-base version, while compact support preserves intersections outside its support. The source explicitly warns that its intended workshop question may have been different.
\subsection{Sources}
\begin{itemize}
\item Local source record: catalogue entry, catalogue entry; canonical \texttt{aim-geometry-notes.json}, \texttt{notes[398]}, Question 1.1. Title: ``Zero-section intersections in a 1-jet space.''
\item AIM 2003 workshop ``Holomorphic curves in contact geometry,'' supplied reference \url{https://aimath.org/WWN/contactgeom2/contactgeom2.pdf}. Original attributions: Tobias Ekholm for an approximately recorded question and Eliashberg for the generating-function observation. This draft uses the supplied text without accessing a PDF.
\item Supplied research summary invokes Chekanov generating-family persistence and quadratic-at-infinity Morse theory without full bibliographic entries. Attempt 1 supplies the closed/open/supported trichotomy and explicit strict contact isotopy.
\item Registry dated 14 August 2026: ``partially\_solved,'' ``PARTIAL-PROGRESS,'' statement status ``exact''; attempt partial, rigor ``proved,'' novelty confidence low. Exact wording does not remove the source author's explicit mathematical uncertainty. The machine-generated, non-peer-reviewed aid requires independent verification before citation.
\end{itemize}
\end{document}
